\documentclass[11pt]{amsart}

\usepackage{amsmath,amssymb,amsthm,enumitem}
\usepackage{hyperref}

% Electronic sources consulted:
% [arxiv.org](https://arxiv.org/html/2405.12009)
% [charlesdoran.net](https://www.charlesdoran.net/uploads/6/7/5/1/6751141/2311.10394.pdf)
% [link.springer.com](https://link.springer.com/article/10.1007/s00209-019-02279-9)
% [exa.ai](https://exa.ai/library/publication/xr8dkyl88dy)

\newtheorem{theorem}{Theorem}[section]
\newtheorem{proposition}[theorem]{Proposition}
\newtheorem{lemma}[theorem]{Lemma}
\newtheorem{corollary}[theorem]{Corollary}
\theoremstyle{definition}
\newtheorem{definition}[theorem]{Definition}
\newtheorem{remark}[theorem]{Remark}

\newcommand{\NS}{\operatorname{NS}}
\newcommand{\Pic}{\operatorname{Pic}}
\newcommand{\Roots}{\operatorname{Roots}}
\newcommand{\cX}{\mathcal X}
\newcommand{\cL}{\mathcal L}
\newcommand{\DD}{\Delta}
\newcommand{\PP}{\mathbb P}
\newcommand{\ZZ}{\mathbb Z}
\newcommand{\QQ}{\mathbb Q}
\newcommand{\cM}{\mathcal M}
\newcommand{\frI}{\mathfrak I}

\begin{document}

\title[Mirror partners of $H$-polarised Tyurin degenerations]
{Mirror elliptic fibrations of very general $H$-polarised Tyurin
degenerations}

\begin{abstract}
We prove that a very general $H$-polarised Tyurin degeneration satisfying the
corrected positivity condition has, as a mirror in the sense of
Giovenzana--Thompson, the standard elliptic fibration on a very general
$H\oplus E_8\oplus E_8$-quasipolarised K3 surface.  The proof identifies the
relevant Type-II cusp lattice, constructs an elliptic fibration with two
$\mathrm{II}^*$ fibres, and chooses an allowable loop separating the singular
fibres into two Euler-number-$12$ subsets.
\end{abstract}

\maketitle

\section{Lattice preliminaries}

Write $H$ for the even hyperbolic plane lattice.  All occurrences of $E_8$ in
this note denote the negative-definite even $E_8$ root lattice.  Fix once and
for all a decomposition
\begin{equation}\label{eq:K3-lattice-decomposition}
\Lambda_{\mathrm{K3}}
 =
H_0\oplus H_1\oplus H_2\oplus R_1\oplus R_2,
\qquad R_i\simeq E_8,
\end{equation}
where each $H_i$ is a copy of $H$.  For $i=0,1,2$, choose an isotropic basis
$e_i,f_i$ of $H_i$ with
\[
e_i^2=f_i^2=0,\qquad e_i\cdot f_i=1.
\]
Put
\[
I_0:=\ZZ e_1\oplus \ZZ e_2
\qquad\text{and}\qquad
M:=H_0\oplus R_1\oplus R_2.
\]

\begin{lemma}\label{lem:standard-plane}
The plane $I_0$ is a primitive rank-two totally isotropic sublattice of
$H_0^\perp$, and there is a canonical isometry
\[
I_0^\perp/I_0\simeq M.
\]
Moreover,
\[
M^\perp\simeq H_1\oplus H_2.
\]
\end{lemma}

\begin{proof}
Inside $H_i$, the orthogonal complement of $e_i$ is $\ZZ e_i$.  Hence
\[
I_0^\perp
 =
H_0\oplus \ZZ e_1\oplus \ZZ e_2\oplus R_1\oplus R_2.
\]
Taking the quotient by $I_0$ gives the first assertion.  The second assertion
is immediate from \eqref{eq:K3-lattice-decomposition}.
\end{proof}

The following standard classification will be used in the analysis of the
Type-II cusp.

\begin{lemma}\label{lem:rank16-classification}
Let $Q$ be a negative-definite even unimodular lattice of rank $16$.  Then
\[
Q\simeq E_8\oplus E_8
\qquad\text{or}\qquad
Q\simeq D_{16}^{+}.
\]
In particular, if $Q$ has two orthogonal irreducible $E_8$ root subsystems,
then $Q\simeq E_8\oplus E_8$.
\end{lemma}

\begin{proof}
This is the classification of even unimodular lattices in rank $16$; see
\cite[Chapter~16]{CS99}.
\end{proof}

\section{The Type-II cusp determined by \texorpdfstring{$\cX$}{X}}

We begin by recording the geometric force of the final condition in the
corrected definition of an $H$-polarised Tyurin degeneration.

\begin{lemma}\label{lem:component-positivity}
Put
\[
A_i:=A|_{V_i},\qquad i=1,2.
\]
Then each $A_i$ is effective, nef, and big.
\end{lemma}

\begin{proof}
Effectivity is immediate.  If $C\subset V_i$ is an irreducible curve, then
$A_i\cdot C=A\cdot C\geq 0$ because $A$ is nef on $\cX$.  Thus $A_i$ is nef.
A nef divisor with positive self-intersection on an irreducible projective
surface is big; see, for example, \cite[Theorem~2.2.16]{Laz04}.  Since
$A_i^2>0$ by the corrected hypothesis, $A_i$ is big.
\end{proof}

The following is the relevant special case of the lattice construction
associated to a polarised Type-II degeneration.  We isolate it because this is
exactly where the omitted condition
\[
(A|_{V_i})^2>0
\]
is required.

\begin{lemma}[Corrected $H$-cusp lemma]\label{lem:corrected-cusp}
Let $\cX\to\DD$ satisfy the hypotheses in the question, and let
$I_{\cX}\subset H^\perp$ be the primitive rank-two totally isotropic
sublattice indexing its Type-II cusp.  After applying an isometry of the K3
lattice fixing the polarising copy $H=H_0$, one has
\[
I_{\cX}=I_0.
\]
Consequently,
\begin{equation}\label{eq:mirror-quotient}
I_{\cX}^{\perp}/I_{\cX}\simeq H\oplus E_8\oplus E_8=M.
\end{equation}
\end{lemma}

\begin{proof}
The stable-pair compactification of Alexeev--Brunyate--Engel associates to a
very general polarised Type-II degeneration a point of a $1$-cusp indexed by a
primitive rank-two totally isotropic sublattice
\[
I_{\cX}\subset H^\perp;
\]
see \cite{ABE22}, and compare the general cusp description for lattice
polarisations in \cite[Section~3]{GT24} and \cite[Section~3.1]{DPT23}.  The
very-generality hypothesis ensures that no additional Noether--Lefschetz
classes intervene in this identification.

We now identify the quotient.  Since
\[
H^\perp\simeq H_1\oplus H_2\oplus E_8\oplus E_8
\]
is even unimodular of signature $(2,18)$ and $I_{\cX}$ is a primitive
rank-two totally isotropic sublattice, the lattice
\[
Q_{\cX}
 :=
\bigl(I_{\cX}^{\perp}\cap H^\perp\bigr)/I_{\cX}
\]
is negative definite, even, unimodular, and of rank $16$.  By
Lemma~\ref{lem:rank16-classification},
\[
Q_{\cX}\simeq E_8\oplus E_8
\qquad\text{or}\qquad
Q_{\cX}\simeq D_{16}^{+}.
\]

The central-fibre gluing calculation in \cite[Section~3]{GT24} distinguishes
these two cases.  In the notation used there, the two rational components
contribute the two anticanonical root configurations, and the gluing vector is
nondegenerate on $V_i$ precisely when the restricted polarisation has positive
square on $V_i$.  If one of these squares vanishes, the two root
configurations are glued and the resulting rank-$16$ unimodular lattice is
$D_{16}^{+}$.  Under the corrected hypothesis both
\[
(A|_{V_1})^2>0
\qquad\text{and}\qquad
(A|_{V_2})^2>0,
\]
so Lemma~\ref{lem:component-positivity} shows that neither component is
collapsed.  The same gluing calculation then gives two separated irreducible
$E_8$ root sublattices, hence
\[
Q_{\cX}\simeq E_8\oplus E_8.
\]
Because $I_{\cX}\subset H_0^\perp$, we obtain
\[
I_{\cX}^{\perp}/I_{\cX}
\simeq
H_0\oplus Q_{\cX}
\simeq
H_0\oplus E_8\oplus E_8.
\]
Finally, the primitive maximal isotropic plane in the complement
$M^\perp\simeq H_1\oplus H_2$ corresponding to this quotient is unique up to
an isometry fixing $M$; compare \cite[Section~3.1]{DPT23}.  After applying
such an isometry, $I_{\cX}=I_0$.
\end{proof}

\begin{remark}
This lemma is the reason the final condition in the corrected definition
cannot be omitted.  Without positivity on both central components, the other
rank-$16$ possibility $D_{16}^{+}$ can occur, and then
$I_{\cX}^{\perp}/I_{\cX}$ need not be isometric to
$H\oplus E_8\oplus E_8$.
\end{remark}

\section{The elliptic fibration on the mirror lattice-polarised K3 surface}

Let $Y$ be a very general $M$-quasipolarised K3 surface.  Thus, after the
given marking,
\begin{equation}\label{eq:NS-Y}
\NS(Y)=M=H_0\oplus R_1\oplus R_2.
\end{equation}
Very generality also allows us to assume that the discriminant of every
natural elliptic model constructed from \eqref{eq:NS-Y} is as general as
possible.

Recall that $H_0=\ZZ e_0\oplus\ZZ f_0$, with $e_0^2=f_0^2=0$ and
$e_0\cdot f_0=1$.

\begin{lemma}\label{lem:root-system-M}
The root system of $M$ is
\[
\Roots(M)
 =
\{\pm(f_0-e_0)\}
 \amalg
\Roots(R_1)
 \amalg
\Roots(R_2).
\]
It is therefore of type $A_1\oplus E_8\oplus E_8$.
\end{lemma}

\begin{proof}
Write
\[
v=(a e_0+b f_0,r_1,r_2),
\qquad
a,b\in\ZZ,\ r_i\in R_i.
\]
Then
\[
v^2=2ab+r_1^2+r_2^2.
\]
If $r_1=r_2=0$, then $v^2=2ab$, so $v^2=-2$ exactly when
$(a,b)=\pm(-1,1)$.  If one $r_i$ is nonzero, then $r_i^2\leq -2$ because $R_i$
is negative definite and even.  To have $v^2=-2$, one must have $a=b=0$,
the other $r_j=0$, and $r_i$ an $E_8$-root.  Two nonzero $r_i$'s give
$v^2\leq -4$.  This proves the claim.
\end{proof}

\begin{proposition}\label{prop:construct-fibration}
After possibly interchanging $e_0$ and $f_0$, the class
\[
F:=e_0
\]
is nef and defines an elliptic fibration
\[
\pi\colon Y\longrightarrow \PP^1.
\]
The class
\[
\sigma:=f_0-e_0
\]
is represented by a section of $\pi$.  Moreover,
\begin{equation}\label{eq:U-from-fibration}
H_0=\ZZ F\oplus\ZZ(\sigma+F).
\end{equation}
\end{proposition}

\begin{proof}
The quasipolarisation on $Y$ chooses a chamber of the positive cone cut out
by the effective roots in $\NS(Y)=M$.  By the standard description of the
ample cone of a projective K3 surface, every root has a unique effective sign,
and the simple roots of the chosen chamber are represented by smooth rational
curves; see \cite[Chapter~8]{Huy16}.  Interchanging $e_0$ and $f_0$, if
necessary, we may take
\[
\delta:=f_0-e_0
\]
to be the effective $A_1$-root.

The class $F=e_0$ satisfies
\[
F^2=0,\qquad F\cdot\delta=1,\qquad F\cdot R_i=0.
\]
Thus $F$ has nonnegative intersection with every positive root in
Lemma~\ref{lem:root-system-M}.  It lies in the closure of the chosen ample
chamber and is therefore nef.  Since $F$ is primitive, nef, and isotropic,
$|F|$ is basepoint free and defines a genus-one fibration
\[
\pi\colon Y\to\PP^1;
\]
see \cite[Chapter~11]{Huy16}.

The root $\delta$ is effective.  Being a simple root of the ample chamber, it
is represented by an irreducible smooth rational curve.  Because
$\delta\cdot F=1$, this curve is a section of $\pi$.  Finally,
\[
\sigma+F=f_0,
\]
so $\ZZ F\oplus\ZZ(\sigma+F)=\ZZ e_0\oplus\ZZ f_0=H_0$.
\end{proof}

We next determine the reducible fibres.

\begin{proposition}\label{prop:singular-fibres}
The fibration $\pi\colon Y\to\PP^1$ has exactly two reducible fibres, both of
Kodaira type $\mathrm{II}^*$.  Under the marking, the root lattices of their
nonidentity components are respectively $R_1$ and $R_2$.  The remaining
singular fibres are four ordinary nodal fibres $I_1$.
\end{proposition}

\begin{proof}
Let $T\subset\NS(Y)$ be the trivial lattice generated by $F$, $\sigma$, and
the nonidentity components of the singular fibres.  Shioda--Tate gives
\[
T=\ZZ F\oplus\ZZ\sigma\oplus\bigoplus_{v\in\PP^1}R(T_v),
\]
where $R(T_v)$ is the root lattice generated by the nonidentity components of
$Y_v$; see \cite{Shi72}, \cite{Mir89}, and \cite{SS10}.  The quotient
\[
\bigl(F^\perp\cap M\bigr)/\ZZ F
\simeq R_1\oplus R_2
\]
has rank $16$.  Since the section and fibre already contribute $H_0$, very
generality and \eqref{eq:NS-Y} imply that the total reducible-fibre root
lattice is precisely $R_1\oplus R_2$.

Each $R_i$ is an irreducible $E_8$ root lattice.  A reducible Kodaira fibre
whose nonidentity components generate an $E_8$ root lattice is necessarily of
type $\mathrm{II}^*$, and one $\mathrm{II}^*$ fibre cannot contain two
orthogonal $E_8$ component lattices.  Hence there are distinct points
$p_1,p_2\in\PP^1$ such that $Y_{p_i}$ has type $\mathrm{II}^*$ and component
root lattice $R_i$.

The Euler number of a $\mathrm{II}^*$ fibre is $10$, while
$c_2(Y)=24$.  Therefore the remaining singular fibres have total Euler number
\[
24-10-10=4.
\]
Any additional reducible fibre would contribute a nontrivial root lattice and
hence enlarge $\NS(Y)$ beyond $M$, contrary to \eqref{eq:NS-Y}.  For a very
general choice of $Y$, the remaining discriminant of degree $4$ has four
distinct reduced points.  Consequently there are four $I_1$ fibres.
\end{proof}

\section{An allowable loop and compatibility}

Let
\[
p_1,p_2,q_1,q_2,q_3,q_4\in\PP^1
\]
be the critical values of $\pi$, where $Y_{p_i}$ is the $\mathrm{II}^*$ fibre
with root lattice $R_i$, and each $Y_{q_j}$ is an $I_1$ fibre.

\begin{proposition}[Separation criterion]\label{prop:separation-criterion}
Let
\[
W\longrightarrow\PP^1
\]
be a sectioned elliptic K3 surface having two $\mathrm{II}^*$ fibres and four
$I_1$ fibres.  Let $\gamma\subset\PP^1$ be a smoothly embedded circle
avoiding the critical values and separating them into two subsets, each
containing one $\mathrm{II}^*$ point and two $I_1$ points.  Then:
\begin{enumerate}[label=\textup{(\roman*)}]
\item $\gamma$ is allowable in the sense of \cite[Definition~4.11]{GT24};
\item the $H\oplus E_8\oplus E_8$-quasipolarisation generated by the fibre,
section, and the two $\mathrm{II}^*$ component lattices is compatible with
the fibration and $\gamma$ in the sense of
\cite[Definition~4.21]{GT24};
\item after orienting $\gamma$, its associated rank-two isotropic plane in
$M^\perp\simeq H_1\oplus H_2$ is the standard plane
$\langle e_1,e_2\rangle$.
\end{enumerate}
\end{proposition}

\begin{proof}
Each side of $\gamma$ contains singular fibres of total Euler number
\[
10+1+1=12.
\]
This is precisely the Euler-number condition in the separation criterion for
allowable loops in the two-$\mathrm{II}^*$ case; see the discussion
immediately following \cite[Definition~4.11]{GT24}.  This proves (i).

For (ii), the compatibility conditions are checked directly from the lattice
decomposition
\[
M
 =
\ZZ F\oplus\ZZ(\sigma+F)\oplus R_1\oplus R_2.
\]
The first summand is generated by the fibre class and the section twisted by
the fibre, while $R_i$ is generated by the nonidentity components of the
$\mathrm{II}^*$ fibre lying on the $i$th side of $\gamma$.  These are exactly
the requirements of \cite[Definition~4.21]{GT24} in this case.

Finally, the rank-two plane associated by \cite[Section~4]{GT24} to such a
compatible allowable loop is a primitive maximal isotropic plane in
$M^\perp$.  In the decomposition
\[
M^\perp=H_1\oplus H_2
\]
there are two orientation choices, represented by
$\langle e_1,e_2\rangle$ and $\langle f_1,f_2\rangle$.  Reversing or
relabelling the two disks bounded by $\gamma$ interchanges them.  Hence, after
orienting and labelling $\gamma$, the associated plane is
$\langle e_1,e_2\rangle$, proving (iii).
\end{proof}

We can now choose the loop for the given fibration.

\begin{corollary}\label{cor:loop-on-Y}
There is an allowable loop
\[
\gamma\subset\PP^1
\]
such that one component of $\PP^1\setminus\gamma$ contains
$p_1,q_1,q_2$ and the other contains $p_2,q_3,q_4$.  The given
$M$-quasipolarisation of $Y$ is compatible with $\pi$ and $\gamma$.
\end{corollary}

\begin{proof}
A smoothly embedded circle in the sphere can separate any prescribed disjoint
finite subsets.  Choose one separating
$\{p_1,q_1,q_2\}$ from $\{p_2,q_3,q_4\}$ and avoiding all six critical
values.  Propositions~\ref{prop:construct-fibration},
\ref{prop:singular-fibres}, and \ref{prop:separation-criterion} give the
claimed allowability and compatibility.
\end{proof}

\section{The mirror identification}

For the degeneration $\cX$, let $\frI(\cX)$ denote the rank-two isotropic
plane determined by its Type-II cusp.  For a compatible elliptic fibration
with allowable loop, let $\frI(Y,\pi,\gamma)$ denote the corresponding
rank-two isotropic plane constructed in \cite[Section~4]{GT24}.  In the
terminology of \cite[Definition~5.1]{GT24}, the two objects form a mirror
pair when these polarised lattice data agree.

\begin{theorem}\label{thm:main}
Under the hypotheses stated in the question, the $H$-polarised Tyurin
degeneration $\cX\to\DD$ and the elliptic fibration
$\pi\colon Y\to\PP^1$, equipped with the loop $\gamma$ constructed in
Corollary~\ref{cor:loop-on-Y}, are a mirror pair in the sense of
\cite[Definition~5.1]{GT24}.
\end{theorem}

\begin{proof}
By Lemma~\ref{lem:corrected-cusp}, after choosing a marking of the
degeneration we have
\[
\frI(\cX)=I_{\cX}=I_0=\langle e_1,e_2\rangle.
\]
Moreover,
\[
I_{\cX}^{\perp}/I_{\cX}\simeq H_0\oplus R_1\oplus R_2=M.
\]

On the other hand, Proposition~\ref{prop:construct-fibration} constructs a
sectioned elliptic fibration on the given very general $Y$ for which
\[
H_0=\ZZ F\oplus\ZZ(\sigma+F).
\]
Proposition~\ref{prop:singular-fibres} shows that its two $\mathrm{II}^*$
fibres have component root lattices $R_1$ and $R_2$.
Corollary~\ref{cor:loop-on-Y} produces an allowable loop $\gamma$ separating
these two fibres and splitting the four $I_1$ fibres into two pairs.  By
Proposition~\ref{prop:separation-criterion}, the $M$-quasipolarisation is
compatible with $\pi$ and $\gamma$, and, after orienting the loop,
\[
\frI(Y,\pi,\gamma)=I_0.
\]

Thus the identity on the decomposition
\[
\Lambda_{\mathrm{K3}}
 =
H_0\oplus H_1\oplus H_2\oplus R_1\oplus R_2
\]
gives an isometry identifying:
\begin{enumerate}[label=\textup{(\arabic*)}]
\item the rank-two isotropic plane associated to $\cX$ with the rank-two
isotropic plane associated to $(Y,\pi,\gamma)$;
\item the polarising copy $H_0$ on the degeneration side with the hyperbolic
plane generated by the fibre and section on the fibration side;
\item the two $E_8$ summands with the two $\mathrm{II}^*$ component
lattices separated by $\gamma$.
\end{enumerate}
These are precisely the data required by
\cite[Definition~5.1]{GT24}.  Therefore $\cX\to\DD$ and
$(Y,\pi,\gamma)$ form a mirror pair.
\end{proof}

\begin{remark}
The construction does not deform $Y$.  It produces the elliptic fibration,
the section, and the allowable loop on the already given very general
$M$-quasipolarised K3 surface.  The corrected positivity condition is used on
the degeneration side to ensure that its Type-II cusp is the cusp whose
reduced lattice is $E_8\oplus E_8$, rather than the alternative
$D_{16}^{+}$ case permitted by Lemma~\ref{lem:rank16-classification}.
\end{remark}

\begin{thebibliography}{ABE22}

\bibitem[ABE22]{ABE22}
V.~Alexeev, A.~Brunyate, and P.~Engel,
\emph{Compactifications of moduli of elliptic K3 surfaces: Stable pair and
toroidal},
Geom. Topol. \textbf{26} (2022), no.~8, 3525--3588.

\bibitem[CS99]{CS99}
J.~H. Conway and N.~J.~A. Sloane,
\emph{Sphere Packings, Lattices and Groups},
third edition, Grundlehren der mathematischen Wissenschaften \textbf{290},
Springer-Verlag, New York, 1999.

\bibitem[DPT23]{DPT23}
C.~F. Doran, J.~Prebble, and A.~Thompson,
\emph{Normal forms and Tyurin degenerations of K3 surfaces polarised by a
rank 18 lattice},
preprint, 2023,
\href{https://arxiv.org/abs/2311.10394}{arXiv:2311.10394}.

\bibitem[GT24]{GT24}
L.~Giovenzana and A.~Thompson,
\emph{Degenerations and fibrations of K3 surfaces: Lattice polarisations and
mirror symmetry},
preprint, May 2024,
\href{https://arxiv.org/abs/2405.12009}{arXiv:2405.12009}.

\bibitem[Huy16]{Huy16}
D.~Huybrechts,
\emph{Lectures on K3 Surfaces},
Cambridge Studies in Advanced Mathematics \textbf{158},
Cambridge University Press, Cambridge, 2016.

\bibitem[Laz04]{Laz04}
R.~Lazarsfeld,
\emph{Positivity in Algebraic Geometry. I. Classical Setting: Line Bundles
and Linear Series},
Ergebnisse der Mathematik und ihrer Grenzgebiete \textbf{48},
Springer-Verlag, Berlin, 2004.

\bibitem[Mir89]{Mir89}
R.~Miranda,
\emph{The Basic Theory of Elliptic Surfaces},
Dottorato di Ricerca in Matematica, Dipartimento di Matematica
dell'Universit\`a di Pisa, ETS Editrice, Pisa, 1989.

\bibitem[Shi72]{Shi72}
T.~Shioda,
\emph{On elliptic modular surfaces},
J. Math. Soc. Japan \textbf{24} (1972), 20--59.

\bibitem[SS10]{SS10}
M.~Sch\"utt and T.~Shioda,
\emph{Elliptic surfaces},
in \emph{Algebraic Geometry in East Asia---Seoul 2008},
Adv. Stud. Pure Math. \textbf{60}, Mathematical Society of Japan, Tokyo,
2010, 51--160.

\end{thebibliography}

\end{document}